% ============================================================================
% Appendix C: Derivation of Einstein Equations
% ============================================================================

\section{Derivation of the Einstein equations from $\Gam$}
\label{app:einstein}

This appendix presents the derivation chain from the coherence matrix $\Gam$ to the
Einstein field equations. The derivation proceeds in five steps: (1)~commutativity of the
macroscopic algebra; (2)~emergent temporal manifold; (3)~emergent spatial manifold,
computed as a Gelfand spectrum; (4)~product spectral triple $M^4 = \mathbb{R} \times S^3$;
(5)~Einstein--Hilbert action from the spectral action, with Lovelock's theorem as a
supplementary argument. The mathematics uses no axioms beyond A1--A5; the reading of
$M^4$ as physical spacetime is an interpretation \status{I}, and the Standard-Model part
of the spectral action is imported from the Connes--Chamseddine finite triple rather than
derived (Step~5).

\paragraph{Retractions in this appendix (2026-09-25).} An earlier version of this appendix
rested on three constructions that are retracted \status{$\times$}: the axis-labelled
sector decomposition $7 = 1_O \oplus \mathbf{3}_{A,S,D} \oplus \bar{\mathbf{3}}_{L,E,U}$
(row 48a: $SU(3) = \mathrm{Stab}_{G_2}(e_O)$ acts irreducibly on the six non-$O$ axes, and
no three of them span an invariant subspace; the complexified
$\mathbb{C}^7 = \mathbb{C} \oplus \mathbf{3} \oplus \bar{\mathbf{3}}$ is classical
representation theory, G\"unaydin--G\"ursey 1973, and gives colour, not space); a real
structure of KO-dimension~6 on $H_{\mathrm{int}} = \mathbb{C}^7$ (none exists: with
$J^2 = 1$, $J\chi = -\chi J$ forces the $\chi = \pm 1$ eigenspaces to have equal
dimension, impossible in odd dimension); and the summed clock of $M$ holons with
$7^M$ readings (it has $6M+1$ readings and the fixed period $2\pi/\omega_0$). The steps
below are restated without them.

\subsection{Step 1: Commutativity of the macroscopic algebra (T-117)}

Consider a composite system of $M$ holons with total Hilbert space
$\Hil_M = \bigotimes_{m=1}^M \Hil_{\mathrm{int}}^{(m)}$, $\dim(\Hil_M) = 7^M$, and
observable algebra $A_M = \bigotimes_m A_{\mathrm{int}}^{(m)}$ where
$A_{\mathrm{int}} = \mathbb{C} \oplus M_3(\mathbb{C}) \oplus M_3(\mathbb{C})$
(T-53).

Define macroscopic averages over a region $\Lambda_\ell(x)$ containing
$|\Lambda_\ell(x)|$ holons:
\begin{equation}
  \bar{O}(x) := \frac{1}{|\Lambda_\ell(x)|} \sum_{m \in \Lambda_\ell(x)} O^{(m)}
  \label{eq:macro-average}
\end{equation}

\begin{theorembox}[T-117: Commutativity \status{T}]
For the composite system satisfying the primitivity condition (T-39a) and finite-range
Gap coupling, the macroscopic observables commute at spacelike separation in the
thermodynamic limit:
\begin{equation}
  [\bar{O}_1(x),\, \bar{O}_2(y)] \to 0
  \qquad\text{as } M \to \infty,\; |x-y| > \ell
  \label{eq:commutativity}
\end{equation}
\end{theorembox}

\paragraph{Proof.}
At the microscopic level, $A_M$ is \emph{non-commutative} (it contains matrix algebras
$M_3(\mathbb{C})$). However, the Goderis--Verbeure--Vets quantum central limit theorem~\cite{goderis1989}
(1989) establishes that for a quantum spin system with: (a)~finite interaction range,
(b)~exponential decay of correlations (clustering), the macroscopic averages become
commutative in the thermodynamic limit.

Condition (a) is satisfied by the finite range of the Gap coupling
($\mathrm{Gap} \in [0,1]$, compactness of $(S^1)^{21}$). Condition (b) follows from the
primitivity of $\Lind_0$ (T-39a~\status{T}), which guarantees a unique stationary state
$I/7$ for the linear part and exponential convergence with rate $\Delta(\Lind_0)$. The
exponential decay of correlations follows from the spectral gap. $\square$

\paragraph{Significance.}
Commutativity of the macroscopic algebra is the prerequisite for applying Gelfand's theorem:
a commutative C$^*$-algebra is isomorphic to the algebra of continuous functions on a
compact Hausdorff space. The proof holds for any local observables (an earlier restriction
to a ``$\mathbf{3}+1$-effective sector'' came from the retracted row 48a and is not used).
The restated T-119 and T-120 below do not need T-117: they estimate the commutators of the
macroscopic algebras directly.

\subsection{Step 2: Emergent temporal manifold (T-118)}

\begin{theorembox}[T-118: Temporal manifold \status{T}]
The temporal C$^*$-algebra is $C_0(\mathbb{R})$, obtained as the scaling limit of the
reading algebras $\mathbb{C}^{N+1}$ of the depth register:
\begin{equation}
  \mathbb{C}^{N+1}
  \xrightarrow{\;\Delta t \to 0,\ \text{window} \to \infty\;}
  C_0(\mathbb{R})
  \label{eq:temporal-algebra}
\end{equation}
(with the origin at the first reading the limit is $C_0([0,\infty))$).
\end{theorembox}

\paragraph{Proof (sketch).}
The depth register records the stratal depth $n \in \{0,\dots,N\}$ as an ordered chain of
readings, stored positionally in the O-registers of $\lceil\log_7(N+1)\rceil$ holons under
a Feynman--Kitaev constraint. One state-independent constraint gives the conditional states
exactly $e^{n\Delta t\,\Lind}\rho_0$ at every reading (T-53b~\status{T}). As the chronon
$\Delta t \to 0$ and the window grows without bound on both sides of the origin, the
reading sets $\Delta t\cdot\{0,\dots,N\}$ converge to $\mathbb{R}$ in the pointed Hausdorff
sense, and sampling embeds $C_0(\mathbb{R})$ isometrically into
$\prod_N \mathbb{C}^{N+1}/\bigoplus_N \mathbb{C}^{N+1}$. $\square$

\noindent\emph{Retracted derivation.} An earlier proof took the clock algebra of $M$
holons to be $\mathbb{C}[\mathbb{Z}_{7^M}]$ with effective period $7^M$ and passed to
$C(S^1)$ and then $C_0(\mathbb{R})$ by decompactification. It is retracted
\status{$\times$}: the summed clock of $M$ holons has $6M+1$ readings and the fixed period
$2\pi/\omega_0$; the $7^M$ ordered readings belong to the same registers read
positionally, which is the depth register above. Composite clocks with incommensurate
frequencies give $C(\mathbb{T}^M)$, not $C_0(\mathbb{R})$.

\subsection{Step 3: Emergent spatial manifold (T-119)}

\begin{theorembox}[T-119: Spatial manifold \status{T} as mathematics]
Let $H_1, H_2, H_3$ be three commuting rotation charges of the holon --- a maximal torus of
$\mathrm{U}(3) = \mathrm{Stab}_{\mathrm{SO}(7)}(e_O) \cap C(L_{e_O})$ (two colour Cartan
generators and $iL_{e_O}$; $\operatorname{rank}\mathrm{SO}(7) = 3$). Then:
(a)~the joint spectra of their averages over $M$ holons fill the weight octahedron
$\cong B^3$; (b)~the joint spectra of the fluctuations $\sqrt{M}(\bar H - \mu)$ fill
$\mathbb{R}^3$, and the spatial algebra --- the minimal unitization of $C_0(\mathbb{R}^3)$
--- is
\begin{equation}
  A_{\mathrm{space}} \cong C(S^3);
  \label{eq:spatial-algebra}
\end{equation}
(c)~$\Sigma^3 = S^3$ is closed and spin with a unique smooth structure, and its Dirac triple
satisfies all seven of Connes' conditions~\cite{connes2008reconstruction}, first order and
Poincar\'e duality included; (d)~colour-singlet charges give only two dimensions, so this
$\Sigma^3$ has colour-charged coordinates. The reading of $\Sigma^3$ as physical space is
\status{I}.
\end{theorembox}

\paragraph{Remarks.}
The number of emergent coordinates is the \emph{rank} of the sector algebra --- the number of
independent commuting charges --- not the dimension of a representation space:
$\operatorname{rank}\mathfrak{u}(3) = 3$, whereas $\dim\mathfrak{u}(3) = 9$,
$\dim\mathfrak{su}(3) = 8$, $\dim G_2 = 14$. Point~(d) is a Coleman--Mandula obstacle for
the physical reading: the count $3$ is a count of colour and gives no rotation group
commuting with colour. The colour-commuting route to $3+1$ is Theorem~48c: the colour-singlet
part of the spin factor $\mathfrak h_2(\Oct) \cong \mathbb{R}^{1,9}$ is
$\mathfrak h_2(\mathbb{C}_O)$, of dimension~4 and signature $(1,3)$, on which
$SL(2,\mathbb{C}_O)$ acts as $SO^+(1,3)$ commuting with colour --- \status{T} as
mathematics, \status{C at (L)} as physical spacetime, with (L)~$\Leftrightarrow$~(P) by
Theorem~48e(f)--(i).

\noindent\emph{Retracted derivation.} The earlier proof took $\Sigma^3$ from the
$\{A,S,D\}$-sector of the axis-labelled decomposition $7 = 1_O \oplus \mathbf{3} \oplus
\bar{\mathbf{3}}$, obtained $d_{\mathrm{spectral}} = 3$ from a Weyl law
$N(\lambda) \sim C M \lambda^{\dim\mathbf{3}}$ on $\bigotimes_m \mathbb{C}^3$, and asserted
that the macroscopic algebra of that sector satisfies all seven Connes axioms, reality
inherited from a KO-dimension-6 finite triple. All three steps are retracted
\status{$\times$}: the axis sector is not $SU(3)$-invariant (row 48a); $\bigotimes_{m=1}^M
\mathbb{C}^3$ is finite-dimensional, so its eigenvalue count saturates at $3^M$ and has no
power law; the first-order condition is a constraint on the Dirac operator, not a
consequence, and the Poincar\'e-duality check assumed the manifold it was meant to
conclude; and no KO-dimension-6 real structure exists on $\mathbb{C}^7$. The restated
T-119 computes the spectrum instead of reconstructing it.

\subsection{Step 4: Product spectral triple $M^4$ (T-120, T-120b)}

\begin{theorembox}[T-120: Product manifold \status{T} as mathematics]
For every Riemannian metric on $M^4 = \mathbb{R}\times S^3$ the product of the temporal and
spatial triples,
\begin{equation}
  (A, H, D)_{\mathrm{macro}} = (C_0(\mathbb{R}), L^2(\mathbb{R}), i\partial_\tau)
  \times
  (C(S^3), L^2(S^3, S), D_{S^3}),
  \label{eq:product-triple}
\end{equation}
is a spectral triple on the 4-manifold $M^4 = \mathbb{R} \times S^3$; the three macroscopic
algebras commute in the limit by a direct estimate, $\|[F_i, a]\| = O(M^{-1/2})$. The
product carries no real structure, so no first-order condition arises. Its reading as
physical spacetime inherits the \status{I} of T-119.
\end{theorembox}

\begin{theorembox}[T-120b: Vacuum topology]
(i)~$\Sigma^3 \cong S^3$ \status{T} --- now part of T-119 (minimal unitization of the
fluctuation spectrum $\mathbb{R}^3$), independent of the vacuum.
(ii)~Constant curvature $k = +1$, a de~Sitter spatial section
\status{C at the vacuum symmetry}: the step uses the $SU(3)$-invariance of a unique sector
vacuum, now the hypothesis (SV).
\end{theorembox}

\paragraph{Proof of T-120.}
The temporal algebra $C_0(\mathbb{R})$ (T-118) and the spatial algebra $C(S^3)$ (T-119)
commute in the limit, so their tensor product is the full macroscopic algebra:
\begin{equation}
  A_{\mathrm{macro}} = C_0(\mathbb{R}) \otimes C(S^3)
  \cong C_0(\mathbb{R} \times S^3) = C_0(M^4).
\end{equation}
The spectral triple product formula (Connes 1996; Chamseddine--Connes 1997) gives the Dirac
operator $D = D_{M^4} \otimes 1 + \gamma_5 \otimes D_{\mathrm{int}}$ on the product with the
finite triple. $\square$

\emph{Lorentzian signature} (T-53, overall \status{C}). Two logically separate steps.
\textbf{(i) The $(1,3)$-split}~\status{T}: the Page--Wootters clock register (T-87, steps
1--3,~\status{T}) singles out exactly one timelike direction (the PW time \emph{axis} ---
not to be confused with the seven tick states of the clock register), while the spatial
sector is the three-dimensional $S^3$ of T-119; the split is
$(\operatorname{rank}\mathfrak{u}(1)_O, \operatorname{rank}\mathfrak{u}(3)) = (1,3)$.
\textbf{(ii) The relative sign} $g_{aa} < 0$~\status{T at reflection positivity}: it
follows from requiring the PW generator to be bounded below (Osterwalder--Schrader
reflection positivity across the PW slice), not from an ansatz. The Krein spectral triple
with fundamental symmetry $\beta = \gamma^0 \otimes 1$, in which $D$ is Krein-self-adjoint,
realizes the signature $(+1,-1,-1,-1)$, but is a consistency check rather than a
derivation: the Euclidean set $\{\gamma^0, i\gamma^i\}$ is Krein-self-adjoint just as
exactly with $\beta = 1$. (An earlier version assigned to a KO-dimension-6 real structure
$J$, T-53, the role of fixing the even/odd grading; no such structure exists on
$\mathbb{C}^7$, and the step is withdrawn.) Theorem~48c gives a second route to the same
signature without reflection positivity: the sign of $\det$ on $\mathfrak h_2(\mathbb{C}_O)$,
\status{T} as mathematics, \status{C at (L)} as spacetime.

\paragraph{Former proof of T-120b (retracted in part).}
The earlier proof took the vacuum $\rho^* = \vp(\Gam)$ to be invariant under
$\mathrm{SU}(3) \subset G_2$ acting on a spatial $\mathbf{3}$-sector, and obtained constant
positive curvature from $\Lambda_{\mathrm{Gap}} > 0$ via T-186(c). Both inputs are gone: the
spatial sector was the axis triple of row 48a, the vacuum of the $G_2$-invariant
$V_{\mathrm{Gap}}$ keeps no $SU(3)$ sector pattern (T-64, T-331; the sector values are the
hypothesis (SV)), and T-186(c) (``$\Delta F > 0$ unconditionally'') is retracted
\status{$\times$}. Hence part~(i) now comes from T-119 and part~(ii) is conditional.

\subsection{Step 4b: The gauge group}

An earlier version of this step reduced $A_{\mathrm{int}} = \mathbb{C} \oplus
M_3(\mathbb{C}) \oplus M_3(\mathbb{C})$ to the Connes--Chamseddine algebra
$A_F = \mathbb{C} \oplus \mathbb{H} \oplus M_3(\mathbb{C})$ through the KO-dimension-6 real
structure, the Fano--Higgs line $\{A,E,U\}$ and a quaternionic selection on the
$\{E,U\}$ block, and claimed a Morita equivalence $A_{\mathrm{int}} \sim A_F$ with the
Standard-Model group as $\mathrm{Inn}(A_F)$ (T-175a). This is retracted \status{$\times$}:
Morita equivalence preserves the centre, and $Z(A_{\mathrm{int}}) = \mathbb{C}^3$ (real
dimension~6) differs from $Z(A_F) = \mathbb{C} \oplus \mathbb{R} \oplus \mathbb{C}$ (real
dimension~5); the unitary groups differ as well ($U(1)\times U(3)\times U(3)$, dimension~19,
against $U(1)\times SU(2)\times U(3)$, dimension~13); and the real structure does not exist.
Nothing replaces the claim within the spectral-triple route.

The Standard-Model group is obtained instead from the Clifford system of
$\mathbb{C}\otimes\Oct$: $(SU(3)\times SU(2)\times U(1))/\mathbb{Z}_6$ is the normaliser of
colour in $\mathrm{Spin}(9)$, and the electroweak algebra is the centraliser of colour
(T-326); the doublets are chiral (T-327), families are horizontal (T-328), and one complete
generation with right-handed fields follows from the complexified spinor (T-329). These are
\status{T} as mathematics and \status{C at (Cl)} in UHM, where (Cl$_0$) states that fermions
are vectors of the spinor module $\mathcal{S} = \mathbb{C}\otimes\Oct$.

\subsection{Step 5: Einstein--Hilbert action (T-65, T-121)}

\begin{theorembox}[T-65: Spectral action \status{T} for the Einstein--Hilbert term]
The spectral action on the product $(M^4 \times A_{\mathrm{int}})$ yields the
Einstein--Hilbert action,
\begin{equation}
  S = \frac{1}{16\pi G}\int_{M^4} (R_{\mathrm{scalar}} - 2\Lambda)\,\sqrt{-g}\,d^4x
  + S_{\mathrm{SM}},
  \qquad G_N = \frac{3\pi}{7 f_2 \Lambda_{\mathrm{cut}}^2},
  \label{eq:einstein-hilbert}
\end{equation}
where the heat-kernel coefficient $a_2$ sees the internal space only through
$\mathrm{Tr}(1) = 7$. The Standard-Model part $S_{\mathrm{SM}}$ is imported with
Connes' finite Hilbert space $H_F$, not derived (see below).
\end{theorembox}

\begin{theorembox}[T-121: Lovelock closure \status{T}]
Lovelock's uniqueness argument~\cite{lovelock1971} applies to the emergent $M^4$ and is
supplementary to the spectral one.
\end{theorembox}

\paragraph{Proof.}
Lovelock's theorem~\cite{lovelock1971} states that in 4 dimensions, the most general
symmetric, divergence-free $(0,2)$-tensor constructible from the metric $g_{\mu\nu}$ and
its first two derivatives is:
\begin{equation}
  G_{\mu\nu} + \Lambda g_{\mu\nu}
  \label{eq:lovelock}
\end{equation}
where $G_{\mu\nu} = R_{\mu\nu} - \frac{1}{2}Rg_{\mu\nu}$ is the Einstein tensor.
Its premises are met as follows.

\begin{enumerate}[leftmargin=2em]
  \item \textbf{Smooth four-dimensional manifold (discreteness $\to$ continuity).}
        $M^4 = \mathbb{R}\times S^3$ is a smooth 4-manifold (T-120~\status{T} as
        mathematics), one dimension from the depth register (T-118) and three from the
        rank of $\mathfrak{u}(3)$ (T-119).

  \item \textbf{Covariance.} Four-dimensional diffeomorphism covariance is inherited from
        the diffeomorphism invariance of the Chamseddine--Connes spectral action. An
        earlier leg ``$G_2 \to SU(3) \to SO(3) \subset \mathrm{Diff}(M^4)$'' is retracted
        \status{$\times$}: no non-trivial homomorphism $SU(3) \to SO(3)$ exists, and the
        axis triples are not $SU(3)$ sectors.

  \item \textbf{Signature.} Lovelock's theorem works for any signature, so this gap is
        irrelevant to the argument; the Lorentzian signature is Step~4 (T-53, \status{C}).
        An earlier version derived it from the KO-dimension-6 real structure $J$ via
        opposite-signed eigenvalues of $D_{\mathrm{int}}$ on the temporal and spatial
        blocks and the Connes distance formula; that route is withdrawn with the real
        structure, and the ansatz $g_{\mu\mu} = \chi_{\mu\mu}/\|D_\mu\|^2$ is retired.
\end{enumerate}

The spectral action principle~\cite{chamseddine1997} (Connes--Chamseddine) applied to the
product spectral triple $(M^4, A_{\mathrm{int}})$ yields:
\begin{equation}
  \Tr\bigl(f(D^2/\Lambda_{\mathrm{cut}}^2)\bigr)
  = \frac{1}{16\pi G}\int_{M^4} R_{\mathrm{scalar}}\,\sqrt{-g}\,d^4x
  + \text{cosmological constant}
  + \text{Standard Model}
  + O(\Lambda_{\mathrm{cut}}^{-2})
  \label{eq:spectral-action}
\end{equation}
where $f$ is a smooth cutoff function and $\Lambda_{\mathrm{cut}}$ is the energy cutoff; the
value of $G_N$ needs the cutoff convention $f_2$ and the scale $\Lambda_{\mathrm{cut}}$
\status{D}. The product is a spectral triple but not a real one in Connes' sense, and the
Einstein--Hilbert term does not use a real structure.

\emph{Status of the Standard-Model term.} The earlier text derived the Standard-Model
Lagrangian from the internal triple $(A_{\mathrm{int}}, H_{\mathrm{int}}, D_{\mathrm{int}})$
with KO-dimension~6. This is withdrawn \status{$\times$}: $H_{\mathrm{int}} = \mathbb{C}^7$
carries no such real structure and has dimension~7 while one generation needs 32 states,
so the Standard-Model content comes from Connes' imported $H_F$ (the bimodule derivation
T-178 is retracted as a derivation). The import brings its history: under the ``big
desert'' assumption Chamseddine, Connes and Marcolli (2007) obtained a Higgs mass around
170\,GeV, excluded at 95\% C.L.\ by CDF and D0 in 2008; the Higgs was found at 125\,GeV, and
the 2012 rescue adds a real singlet and a fitted parameter.

The variation $\delta S / \delta g^{\mu\nu} = 0$ yields the Einstein field equations:
\begin{equation}
  G_{\mu\nu} + \Lambda g_{\mu\nu} = 8\pi G\, T_{\mu\nu}
  \label{eq:einstein-equations}
\end{equation}
where $T_{\mu\nu}$ is the energy-momentum tensor of the matter fields. $\square$

\subsection{Summary of the derivation chain}

\begin{equation}
  \underbrace{\Gam \in \Dens(\mathbb{C}^7)}_{\text{Axiom}}
  \xrightarrow{\text{T-118, T-119}}
  \underbrace{M^4 = \mathbb{R} \times S^3}_{\text{manifold}}
  \xrightarrow{\text{T-120, T-65}}
  \underbrace{G_{\mu\nu} + \Lambda g_{\mu\nu} = 8\pi G T_{\mu\nu}}_{\text{Einstein}}
  \label{eq:derivation-chain}
\end{equation}

Every step of the gravitational chain uses either previously proven theorems of UHM (T-39a,
T-53b, T-62, T-118--T-120) or standard mathematical results (Gelfand's theorem, the Dirac
triple of $S^3$, quantum CLT, heat-kernel expansion, Lovelock's theorem), and is
\status{T} as mathematics. What is not derived is stated: the reading of $M^4$ as physical
spacetime \status{I} (or, by the colour-commuting route of Theorem~48c, \status{C at (L)}),
the Lorentzian sign at reflection positivity, the cutoff convention that fixes the value of
$G_N$, and the Standard-Model part of the action, which is imported. Within these limits the
Einstein--Hilbert term is a \emph{consequence} of the same axioms that define consciousness
--- because spacetime and experience are two aspects of the same coherence matrix $\Gam$.
